% ===== 原文 pp.112–116（PDF p.122–126） =====
% [原文 p.112]
利用这些定义，我们可以陈述 Neumann 条件如下：
\[
G(u,w)=F(v,w)
\]
对所有 $w\in TY$ 成立。正如所承诺的，这给出 $k$ 个独立条件，每个线性独立的切向量 $w$ 对应一个。
% [Tutor Note: p.112 此处 Neumann 条件写作 $G(u,w)=F(v,w)$，但同页下文“取 $w=u$ 得 $G(u,v)=G(v,u)=F(u,u)=0$”只对互换形式 $G(v,w)=F(u,w)$ 成立；p.117 的完整粘合条件表（$\phi\in Y$、$G(v,w)=F(u,w)$、$G(s,w)=F(s,w)$）与 (9.1) 的变分计算都给出 $G(v,w)=F(u,w)$（其中 $u^I=\partial_0\phi^I$、$v^I=\partial_1\phi^I$）。故 p.112 这处 $u,v$ 写反，属原文笔误；译文照抄原文形式，正确形式见 p.117 的清单]

\begin{definition}[1]
满足上述粘合条件的对 $(u,v)$ 合在一起被称为切向量的一个允许对（allowed pair）。
\end{definition}

写下了好的场粘合条件之后，现在我们来分析拓扑缺陷条件对允许对施加哪些额外约束；这会同时约束子流形 $Y$ 的选取以及其世界体上 2-形式 $F$ 的选取。首先，共形缺陷条件 $T^1{}_0-\hat T^1{}_0=0$ 等价于
\[
G(u,v)=0
\]
对所有允许对 $(u,v)$ 成立。这是 $F$ 反对称性的自动推论，因为若 $(u,v)$ 是允许对，则 $G(u,v)=G(v,u)=F(u,u)=0$，其中我们对允许对的定义取了 $w=u$。因此 $(Y,F)$ 自动定义玻色 $\sigma$-模型的一个共形缺陷。

进一步要求拓扑条件 $T^1{}_1-\hat T^1{}_1=0$，等价于要求
\[
G(u,u)=-G(v,v)
\]
对所有允许对 $(u,v)$ 成立。这个条件结果等价于一个表面上更强的条件，如下所示。

\begin{proposition}[1]
若 $(Y,F)$ 使得对所有允许对 $(u,v)$ 都有 $G(u,u)=-G(v,v)$，则 $v\in TY$，且 $(v,u)$ 也是允许对。
\end{proposition}

（命题 1 的证明以及其他并非显然的证明见附录 D。）
逆命题同样成立：若 $(Y,F)$ 满足本命题的条件，则取 $w=u$ 得 $G(u,u)=F(v,u)=-F(u,v)=-G(v,v)$。于是我们得到：

\begin{proposition}[2]
设 $Y\subseteq X\times\hat X$ 是子流形，$F$ 是其世界体上的闭 2-形式。则 $(Y,F)$ 给出玻色 $\sigma$-模型的拓扑缺陷，当且仅当满足如下条件：

若 $(u,v)$ 是一对向量，使得 $u\in TY$、$v\in e^*TM$，且对所有 $w\in TY$ 有 $G(v,w)=F(u,w)$，则

$v\in TY$ 也成立，且对所有 $w\in TY$ 有 $G(u,w)=F(v,w)$。
\end{proposition}
% [Tutor Note: p.111–112 的 (9.1) 与 Neumann 条件的推导可补出：把 $G=\diag(g,-\hat g)$、$B=\diag(b,-\hat b)$、$\cF=\dd\cA$ 代入 (9.1)，其被积式正是 $-G(v,w)-B(u,w)-\cF(u,w)=-G(v,w)+F(u,w)$（用 $F=-\cF-e^*B$），故对所有允许变分 $w$ 置零即得 $G(v,w)=F(u,w)$。这可作为判定 p.112 那处 $u,v$ 互换的依据]

% [原文 p.113]
\begin{definition}[2]
满足上述定理之条件的对 $(Y,F)$ 称作定义了一个拓扑双膜（topological bibrane）。
\end{definition}
% [Tutor Note: p.113 “the preceding theorem” 指 p.112 的命题 2（本章并无 Theorem）；原文另有 “A pair … are said” 的主谓不一致（应为 “is said”），译文按“对……被称为”处理]

\section{拓扑双膜的几何}

上面我们得到了拓扑双膜 $(Y,F)$ 的刻画，现在把这个条件稍作重新表述，使之具有更易理解的形式。

我们选择给流形 $M=X\times\hat X$ 配上中性号差度规 $G$；因此让我们记录一些关于不定号差空间子流形理论的基本事实。更完整的讨论见 [60]。

\begin{definition}[3]
设 $Y$ 是 $2n$ 维伪黎曼流形 $M$ 的嵌入子流形，$M$ 配有非退化对称度规 $G$。正交子丛 $(TY)^\perp$ 定义为与整个 $TY$ 都 $G$-正交的向量构成的集合：
\[
(TY)^\perp=\{u\in e^*TM:\ G(u,v)=0\ \ \forall v\in TY\}.
\]
\end{definition}

若 $TY$ 是 $k$ 维的，则 $(TY)^\perp$ 是 $(2n-k)$ 维的。伪黎曼子流形理论与黎曼子流形理论的主要区别在于：度规 $G$ 限制到 $Y$ 上可能出现退化方向。因此子丛 $\Delta=TY\cap(TY)^\perp$ 一般是非平凡的。借用辛几何的术语，我们定义如下三类特殊子流形：

\begin{definition}[4]
视丛 $TY$（作为 $e^*TM$ 的子丛）是否包含（或被包含于）其正交丛 $(TY)^\perp$，我们说：

$TY$ 是各向同性的，若 $TY\subseteq(TY)^\perp$；

$TY$ 是余迷向的，若 $TY\supseteq(TY)^\perp$；

$TY$ 是拉格朗日的，若 $TY=(TY)^\perp$。
\end{definition}

因此维数 $k$ 落在如下范围：各向同性子丛为 $0\le k\le n$，余迷向子丛为 $n\le k\le 2n$，而所有拉格朗日子丛都是 $n$ 维的。（只有当 $G$ 的号差为 $(n,n)$ 时拉格朗日子丛才存在。）在命题 1 的证明中已经表明，对应于拓扑双膜的子流形 $Y$ 具有这样的性质：正交于它的向量也切于该子流形；在我们的分类中，它们就是余迷向子流形。因此在下面我们聚焦于余迷向子流形（拉格朗日子流形作为特例包含在内）。

% [原文 p.114]
定义一个适配子流形 $Y$ 的 $e^*TM$ 的特定标架是方便的。这不像黎曼情形那样直截了当，因为规范的分解 $e^*TM=TY\oplus(TY)^\perp$ 不再可用。不过，对余迷向子流形存在 [59] 如下形式的切丛分解
\[
e^*TM=\underbrace{(TY)^\perp\oplus SY}_{TY}\oplus NY,
\]
其中 $SY$ 是 $(TY)^\perp$ 在 $TY$ 内的补屏幕分布，$NY$ 是 $TY$ 在 $e^*TM$ 中的补横截分布。我们有 $\dim (TY)^\perp=\dim NY=2n-k$，$\dim SY=2k-2n$。

利用这个适配标架，我们现在写下拓扑双膜的等价刻画：

\begin{proposition}[3]
对 $(Y,F)$ 定义一个拓扑双膜，当且仅当 $Y\subseteq X\times\hat X$ 是余迷向子流形，满足 $\ker F=(TY)^\perp$（即 $F$ 的退化方向与 $G|_{TY}$ 的退化方向一致），并且此外在 $SY$ 上
\[
\left(\tilde G^{-1}\tilde F\right)^2=+1,
\]
其中 $\tilde G\equiv G|_{SY}$、$\tilde F\equiv F|_{SY}$。
\end{proposition}

这与 [56] 中给出的 (2,2) 超对称 $\sigma$-模型 A-膜边界条件的刻画惊人地相似，只要把靶空间上反对称辛形式 $\Omega$ 的角色换成中性号差的对称度规 $G$。正如 A-膜被要求关于 $\Omega$ 余迷向，拓扑双膜被要求关于 $G$ 余迷向。此外，对 A-膜，商丛 $TY/(TY)^\Omega$ 上配有一个由 $\Omega^{-1}F$ 给出、平方为 $-1$ 的自同态；类似地，对拓扑双膜，商丛 $TY/(TY)^\perp$ 上配有自同态 $G^{-1}F$，其平方为 $+1$。
% [Tutor Note: 命题 3 的 $(\tilde G^{-1}\tilde F)^2=+1$ 与 p.116 的稳定化条件可以互相核验：在填满空间情形 $SY=TY$，自同态 $G^{-1}F$ 平方为 $+1$ 意味着 $G^{-1}F$ 是殆积结构（$R$ 型），而广义度规 $\cR_G$ 作用为 $(u,\xi)\mapsto(G^{-1}\xi,Gu)$，$F$-旋转条件为 $Fu=e^*\xi$；把 $\xi=Gu$ 代入得 $F u=G u$ 即 $G^{-1}Fu=u$，正是 $\cR_G(\tau^FY)=\tau^FY$ 的 $+1$ 特征条件]

我们考虑两类重要的特殊拓扑双膜。当 $F=0$ 时，“$F$ 的退化方向与 $G|_{TY}$ 的退化方向一致”这一条件蕴含 $TY=(TY)^\perp$，即 $Y$ 是 $X\times\hat X$ 的拉格朗日子流形。这又蕴含局部上 $Y$ 是等距映射 $f:X\to\hat X$ 的图像。这类“等距映射图像”型双膜是拉格朗日 A-膜的类比物。

另一类特殊拓扑双膜是“填满空间”的双膜，即 $Y=X\times\hat X$。此时 $(TY)^\perp$ 平凡，屏幕分布 $SY$ 就是整个 $TY$。命题 3 中的条件于是蕴含 $F$ 是一个辛形式，且 $(G^{-1}F)^2=1$。填满空间双膜是填满空间 A-膜的拓扑双膜类比物。

确实，若一个流形 $M$ 带有辛形式 $F$、殆积结构 $R$ 以及中性号差度规 $G$，并满足相容性要求
\[
F(u,v)=G(Ru,v)
\]
% [原文 p.115]
对所有向量场 $u,v$ 成立，且若此外还要求
\[
\dd F=0,
\]
则称 $M$ 为殆 para-Kähler 流形 [62]。另一种说法是，满足上述要求的三元组 $(F,R,G)$ 称为该流形上的殆双拉格朗日结构 [61]，这一术语的灵感来自如下事实：$R$ 的 $+1$ 与 $-1$ 特征丛构成关于辛形式 $F$ 的两个互补拉格朗日子丛。若其中一个子丛可积，我们称该流形为半 para-Kähler；若两个子丛都可积，则称该流形为 para-Kähler。一个特征丛的可积性不必蕴含另一个特征丛的可积性。确实，在 §9.4 中为保证超对称性，我们只要求正特征丛可积。para-Kähler 流形的研究是不定号差几何中一个内容丰富且正在发展的分支；它们曾在与本文无关的物理语境中出现于 [48]。

重新表述拓扑双膜条件的另一种有用方式是用广义几何的语言（综述见 [58]）。在广义几何的框架中，人们只与 $TM\oplus T^*M$ 及其截面打交道（而不是孤立地处理 $TM$ 或 $T^*M$）。特别地，人们谈论 $M$ 的广义切向量，定义为向量对 $(u,\xi)$，其中 $u\in TM$、$\xi\in T^*M$。

普通几何的对象有广义对应物；例如，定义子流形 $Y\subseteq M$ 的广义切子丛为其切丛与余法丛之和（后者视为 $T^*M$ 的子丛）：
\[
\tau Y=TY\oplus N^*Y.
\]
广义几何提供了融入非零 2-形式 $F$ 的规则：

\begin{definition}[5]
子流形 $Y\subseteq M$ 的 $F$-旋转广义切子丛是满足如下条件的广义切向量 $(u,\xi)$ 的集合：
\[
\tau^FY=\{(u,\xi):\ u\in TY,\ \xi\in T^*M,\ Fu=e^*\xi\}.
\]
这里 $e:Y\to M$ 是嵌入映射，$Fu$ 表示把 2-形式 $F$ 与向量 $u$ 缩并所得的 1-形式。
\end{definition}

我们一直称为切向量“允许对”的东西，正是 $\tau^FY$ 的截面（在把向量 $v$ 降指标为 1-形式 $\xi=Gv$ 之后）。

广义几何中复结构的对应物是广义复结构：
% [原文 p.116]
即满足 $\cJ^2=-1$ 的自同态 $\cJ:TM\oplus T^*M\to TM\oplus T^*M$。带辛形式 $\Omega$ 的辛流形带有广义复结构
\[
\cJ_\Omega=\begin{pmatrix}&-\Omega^{-1}\\ \Omega&\end{pmatrix}
\]
（其中 $T^*M$ 上的坐标按对偶基理解），而 A-膜可以被优雅地刻画为那些使 $F$-旋转广义切子丛被 $\cJ_\Omega$ 稳定化的 $(Y,F)$：

\begin{proposition}[4]
（见文献 [58]）$(Y,F)$ 定义一个辛流形的 A-膜，当且仅当 $\cJ_\Omega$ 作用在 $\tau^FY$ 的一个截面上给出 $\tau^FY$ 的另一个截面，即
\[
\cJ_\Omega(\tau^FY)=\tau^FY.
\]
\end{proposition}

受我们把玻色 $\sigma$-模型的拓扑双膜与 (2,2) $\sigma$-模型的 A-膜作比较这一主题的启发（把 $\Omega$ 换成中性号差度规 $G$，并把殆复结构换成殆积结构），让我们定义广义（殆）积结构为满足 $\cR^2=1$ 的自同态 $\cR:TM\oplus T^*M\to TM\oplus T^*M$。确实，编码度规的广义几何对象就是一个广义殆积结构：

\begin{definition}[6]
按照 [58] 中的定义，广义度规\footnote{不同作者对“广义度规”这一名称应赋予哪个对象意见不一；我们遵循 [58] 的约定。}是如下特定的广义殆积结构：
\[
\cR_G=\begin{pmatrix}&\ G^{-1}\\ G&\end{pmatrix}.
\]
\end{definition}

最后，我们可以紧凑地描述拓扑双膜如下。

\begin{proposition}[5]
设 $Y\subseteq M$、$F$ 是 $Y$ 上的闭 2-形式，则 $(Y,F)$ 满足拓扑双膜的条件（如上一节所定义），当且仅当 $\cR_G$ 作用在 $\tau^FY$ 的一个截面上给出 $\tau^FY$ 的另一个截面，即
\[
\cR_G(\tau^FY)=\tau^FY.
\]
\end{proposition}

\section{(0,1) 超对称 \texorpdfstring{$\sigma$}{σ}-模型的拓扑缺陷}

现在进入我们最关心的主题：(0,1) 超对称 $\sigma$-模型的超对称拓扑缺陷。除玻色场 $\phi^i$ 与 $\hat\phi^i$ 之外，现在还有 $\Sigma$ 上的右移费米子 $\psi^i_+$ 以及 $\hat\Sigma$ 上的 $\hat\psi^i_+$。
